\Usection{Unidad 3}
\subsection*{Introducción al lenguaje C.}

\subsection*{Contenido según programa}
Elementos del lenguaje C. Introducción a la sintaxis del lenguaje C.  Primer ejemplo: Hola Mundo. Identificación de los elementos de sintaxis. Uso del compilador. Tipos de datos, tamaño y declaraciones. Constantes. Declaraciones.  Operadores aritméticos, relacionales y lógicos. Cast. Jerarquía de operadores.  Operadores de evaluación (expresiones condicionales). Operadores de asignación.  Precedencia. Preprocesador. Archivos de cabecera. Encabezado stdio.h. Entrada y salida con formato. Funciones básicas de entrada salida: scanf, printf, getch, getchar.

\begin{ejercicioClase}
    \subsubsection{Responda las siguientes preguntas}
    \begin{enumerate}[a)]
      \item Los programas en C normalmente son escritos en un:
      \item ¿Que programa es el encargado de combinar la salida del compilador, con las funciones de bibliotecas para producir el archivo ejecutable?
      \item ¿Cual es el programa que carga en memoria el programa ejecutable, desde el disco?
      \item ¿Como se llama el programa encargado de convertir los programas escritos en lenguajes de alto nivel, al lenguaje 
        de máquina?
      \item ¿Quien es el encargado de colocar un programa en memoria para que pueda ser ejecutado?
      \item ¿Cual es la unidad mínima de información en una computadora?
    \end{enumerate}


    \subsubsection{Memoria}
    Responda:
    \begin{enumerate}[a)]
      \item ¿Qué es la memoria principal de una computadora?
      \item ¿Qué diferencia existe entre memoria RAM y almacenamiento en disco?
      \item ¿Por qué el programa debe cargarse en memoria antes de ejecutarse?
      \item Mencione dos tipos de memoria que conozca.
    \end{enumerate}

\end{ejercicioClase}

\begin{ejercicioClase}
    \subsubsection{Completar los espacios en blanco}
    \begin{enumerate}
      \item Todo programa en C comienza con la ejecución de la función \underspace.
      \item Todos los cuerpos de las funciones comienzan con \underspace y terminan con \underspace
    \end{enumerate}

    \subsubsection{Completar los espacios en blanco}
    \begin{enumerate}
      \item El archivo de cabecera que permite usar \texttt{printf} y \texttt{scanf} es \underspace.
      \item La instrucción que finaliza la ejecución de la función \texttt{main} es \underspace.
    \end{enumerate}

    \subsubsection{Programas}
    Escribir un programa en C que resuelva los siguientes problemas (uno por cada enunciado)
    \begin{itemize}
      \item Definir las variables, ``unavariable'', ``p345'', ``numero''.
      \item Escribir un mensaje en pantalla, solicitando al usuario ingresar un número entero. Se debe imprimir la
        solicitud, en la cual al final debe incluir dos puntos y dejarse un espacio.
      \item Solicitar al usuario ingresar un número entero y almacenarlo en una variable llamada ``num''.
      \item Imprimir en pantalla el mensaje ``Buen día''.
      \item Declarar una variable de tipo \texttt{float} llamada \texttt{precio} y otra de tipo \texttt{int} llamada \texttt{cantidad}.
      \item Mostrar en pantalla el mensaje: \texttt{Ingrese un caracter: } sin salto de línea final.
      \item Leer un carácter del teclado y almacenarlo en la variable \texttt{letra}.
      \item Declarar tres variables enteras y asignarles los valores 10, 20 y 30 respectivamente.
    \end{itemize}
    
    \subsubsection{Ejercicio}
    Realizar, para cada enunciado, un programa en C.
    \begin{itemize}
        \item Asignar la suma de las variables \textbf{a} y \textbf{b} en \textbf{c}.
      \item Leer 3 números enteros desde el teclado, y almacenarlos en las variables \textbf{p}, \textbf{q} y \textbf{r}.
    \end{itemize}

    \subsubsection{Verdadero o falso}
    Justifique cada respuesta.
    \begin{itemize}
      \item La declaración \texttt{float x;} crea una variable para almacenar números enteros.
      \item El operador lógico \texttt{\&\&} devuelve 1 sólo si ambas condiciones son verdaderas.
    \end{itemize}

    \subsubsection{Verdadero o falso}
    Indicar si las siguientes afirmaciones son verdaderas o falsas. Si son falsas, indicar porqué:
    \begin{itemize}
      \item La función \textbf{printf} siempre imprime al comienzo de una nueva línea.
      \item Los comentarios son mostrados en la pantalla cuando el programa se está ejecutando.
      \item La secuencia de escape \textbackslash \textbf{n} en la función \textbf{printf}, provoca un salto de línea.
    \end{itemize}

\end{ejercicioClase}


\begin{ejercicioClase}
    \subsubsection{Ejercicios de operadores}
    Realizar un programa en C para cada enunciado:
    \begin{itemize}
          \item Dado un número entero \( n \), calcular su cuadrado.
          \item Ingresar dos números y mostrar su suma, resta, multiplicación y división.
          \item Leer un número entero e imprimir su antecesor y su sucesor.
    \end{itemize}

    \subsubsection{Ejercicio}
    Escribir un algoritmo para calcular la distancia recorrida (m) por un móvil que se desplaza con velocidad constante (m/s) durante un tiempo (s).
    La velocidad y el tiempo serán ingresadas por el usuario.

    \subsubsection{Ejercicio}
    Escribir un algoritmo para obtener el promedio simple de un estudiante a partir de las tres notas parciales.
    Las notas serán introducidas una a una por el usuario.

    \subsubsection{Ejercicio}
    Dado un número de tres cifras, imprimir la suma de sus dígitos. El usuario ingresa el número.

\end{ejercicioClase}


\subsubsection{Completar los espacios en blanco}
\begin{enumerate}
  \item Todas las declaraciones terminan con \underspace.
  \item La función \underspace de la biblioteca estándar permite mostrar información en la pantalla.
  \item La función \underspace de la biblioteca estándar permite ingresar información desde el teclado.
\end{enumerate}
    
\subsubsection{Completar los espacios en blanco}
\begin{enumerate}
  \item Para definir una constante simbólica se utiliza la directiva del preprocesador \underspace.
  \item El operador aritmético que tiene mayor precedencia que la suma y la resta es \underspace.
  \item Para convertir explícitamente un valor flotante a entero se utiliza el mecanismo llamado \underspace.
\end{enumerate}
    
\subsubsection{Verdadero o falso}
Justifique cada respuesta.
\begin{itemize}
  \item La expresión \texttt{3 + 4 * 2} se evalúa de izquierda a derecha sin considerar precedencia.
  \item La instrucción \texttt{getchar()} permite leer un único carácter del teclado.
  \item La expresión \texttt{(int)3.9} produce el valor 4.
  \item El operador \texttt{==} se utiliza para asignar valores a una variable.
\end{itemize}


\subsubsection{Verdadero o falso}
Indicar si las siguientes afirmaciones son verdaderas o falsas. Si son falsas, indicar porqué:
\begin{itemize}
  \item Todas las variables deben definirse antes de ser utilizadas.
  \item El operador de resto (\%) solamente puede ser utilizado en números enteros.
  \item Las variables \textbf{numero} y \textbf{NUmerO} son idénticas.
  \item Los operadores aritméticos \textbf{/,*,+,-,\%} tienen todos el mismo nivel de precedencia.
\end{itemize}

\subsubsection{Ejercicio}
Dado un número real ingresado por el usuario, convertirlo a entero utilizando \textit{cast} y mostrar ambos valores.

\subsubsection{Ejercicio}
Leer dos números enteros e intercambiar sus valores. Mostrar los valores antes y después del intercambio.

\subsubsection{Ejercicio}
Leer un carácter y mostrar su valor ASCII (utilizar cast a \texttt{int}).

\subsubsection{Ejercicio}
Escribir un programa que lea el radio de un círculo y calcule su área y su perímetro.  
Usar \(\pi = 3.14159\).

\subsubsection{Ejercicio}
Ingresar un número entero y determinar si es múltiplo de 3 y de 5 simultáneamente.

\subsubsection{Ejercicio}
Ingresar un número real y mostrar su valor truncado, redondeado hacia arriba y redondeado hacia abajo (usar cast y funciones del encabezado \texttt{math.h}).

\subsubsection{Ejercicio}
Ingresar dos números reales y calcular su promedio.

\subsubsection{Ejercicio}
Ingresar un número de segundos y convertirlo en horas, minutos y segundos.




\subsubsection{Responda las siguientes preguntas}
\begin{enumerate}[a)]
  \item ¿Qué diferencia existe entre un programa fuente y un programa objeto?
  \item ¿Qué tipo de archivo genera el preprocesador y qué extensiones suele utilizar?
  \item ¿Qué tareas realiza el compilador durante el proceso de traducción?
  \item ¿Por qué es necesario el enlazador para generar un archivo ejecutable?
  \item ¿Qué relación existe entre el sistema operativo y la ejecución de un programa compilado?
  \item ¿Qué componente del sistema operativo asigna memoria y controla la ejecución del programa?
  \item ¿Cuál es la diferencia entre código máquina y lenguaje ensamblador?
\end{enumerate}

\subsubsection{Clasifique}
Indique a qué fase del proceso de construcción de un programa pertenece cada descripción:
\begin{enumerate}[a)]
  \item Reemplaza las directivas \texttt{\#include} y macros. \hfill \rule{4cm}{0.1pt}
  \item Traduce el código fuente a lenguaje máquina. \hfill \rule{4cm}{0.1pt}
  \item Combina módulos y bibliotecas en un ejecutable. \hfill \rule{4cm}{0.1pt}
  \item Genera un archivo intermedio con código ensamblador. \hfill \rule{4cm}{0.1pt}
\end{enumerate}

\subsubsection{Verdadero o falso}
Justifique cada respuesta.
\begin{enumerate}[a)]
  \item El preprocesador genera directamente el archivo ejecutable. \\
        \vspace{0.5cm}
  \item El compilador transforma código C en código máquina. \\
        \vspace{0.5cm}
  \item El enlazador sólo se utiliza cuando un programa tiene más de un archivo. \\
        \vspace{0.5cm}
  \item Un archivo \texttt{.o} puede ejecutarse directamente en el sistema operativo. \\
        \vspace{0.5cm}
\end{enumerate}

\subsubsection{Completar}
\begin{enumerate}[a)]
  \item El archivo ejecutable es producido por el \rule{5cm}{0.1pt}.
  \item La fase encargada de expandir macros y directivas \texttt{\#include} es el \rule{1cm}{0.1pt}.
  \item El archivo fuente en C suele tener extensión \rule{5cm}{0.1pt}.
\end{enumerate}

\subsubsection{Aplicar}
Dado el siguiente código simple:

\begin{verbatim}
#include <stdio.h>

int main(void) {
    printf("Hola\n");
    return 0;
}
\end{verbatim}

Responda:
\begin{enumerate}[a)]
  \item ¿Qué archivo genera cada etapa del proceso (preprocesar, compilar, ensamblar, enlazar)?
  \item Indique qué comando de \texttt{gcc} utilizaría para obtener sólo el archivo preprocesado.
  \item Indique qué comando de \texttt{gcc} utilizaría para generar únicamente el archivo objeto.
\end{enumerate}

\subsubsection{Ordene las fases del proceso de construcción de un programa}
Indique con números del 1 al 4 el orden correcto:
\begin{enumerate}[a)]
  \item \rule{1cm}{0.1pt} Enlazado  
  \item \rule{1cm}{0.1pt} Preprocesado  
  \item \rule{1cm}{0.1pt} Ensamblado  
  \item \rule{1cm}{0.1pt} Compilación  
\end{enumerate}

\subsubsection{Ordenar de forma correcta}
Se desea compilar el programa \texttt{saludo.c} utilizando la línea de comandos (terminal) mediante el programa gcc.
Dada la siguiente lista de pasos, ordene los mismo para producir el archivo ejecutable \texttt{saludo.out}.
Tenga en cuenta que alguno de los comandos pueden ser incorrectos.
Recuerde: \textbf{pre}-procesador, \textbf{com}pilador, \textbf{en}lazador.

\begin{enumerate}
  \item \texttt{gcc -c programa.i}
  \item \texttt{gcc saludo.o -o saludo.c}
  \item \texttt{gcc saludo.c -c saludo}
  \item \texttt{gcc saludo.i -o saludo.out}
  \item \texttt{gcc -E saludo.c > programa.i}
  \item \texttt{gcc -c programa.i}
  \item \texttt{gcc -E programa.c > programa.out}
  \item \texttt{gcc programa.o -o saludo.out}
\end{enumerate}


\subsubsection{Unidad Central de Procesamiento}
Responda:
\begin{enumerate}[a)]
  \item ¿Qué significa la sigla CPU?
  \item ¿Cuáles son las principales partes que componen la CPU?
  \item ¿Qué función cumple la Unidad Aritmética y Lógica (ALU)?
  \item ¿Qué función cumple la Unidad de Control?
\end{enumerate}

\subsubsection{Contador de Programa}
Responda:
\begin{enumerate}[a)]
  \item ¿Qué es el contador de programa (Program Counter)?
  \item ¿Qué información almacena?
  \item ¿Qué sucede con su valor después de ejecutar una instrucción?
\end{enumerate}

\subsubsection{Buses}
Responda las siguientes preguntas:
\begin{enumerate}[a)]
  \item ¿Qué es un bus en una computadora?
  \item ¿Qué tipos de buses existen en una arquitectura básica?
  \item ¿Qué información transporta el bus de datos?
  \item ¿Qué función cumple el bus de direcciones?
\end{enumerate}

\subsubsection{Dispositivos de entrada y salida}
Clasifique los siguientes dispositivos como \textbf{Entrada (E)}, \textbf{Salida (S)} o \textbf{Entrada/Salida (E/S)}:

\begin{enumerate}[a)]
  \item Teclado \hfill \rule{2cm}{0.1pt}
  \item Monitor \hfill \rule{2cm}{0.1pt}
  \item Disco rígido \hfill \rule{2cm}{0.1pt}
  \item Impresora \hfill \rule{2cm}{0.1pt}
  \item Mouse \hfill \rule{2cm}{0.1pt}
\end{enumerate}

\subsubsection{Análisis}
Lea la siguiente situación y responda:

Un estudiante escribe un programa en C, lo compila y luego lo ejecuta desde la terminal.

\begin{enumerate}[a)]
  \item ¿En qué momento interviene el compilador?
  \item ¿Qué componente del sistema operativo carga el programa en memoria?
  \item ¿Qué componente del hardware ejecuta las instrucciones del programa?
\end{enumerate}